% ============================================================================
% Fichier  : front/02_mode_emploi.tex (Mode d'emploi)
% Version  : 1.0 -- novembre 2024
% ============================================================================
% 30 h de cours magistral = 15 séances de 2 h ; TD (fiches, exercices) et TP
% (dix labs) faits par l'étudiant chez lui, hors des 30 h.
% ============================================================================

\chapter*{Mode d'emploi du cours}
\addcontentsline{toc}{chapter}{Mode d'emploi du cours}
\markboth{Mode d'emploi}{Mode d'emploi}

\section*{Structure générale}

Le cours compte un prologue et six parties, soit seize chapitres, pour
30~heures de cours magistral réparties en quinze séances de deux heures.
Les travaux dirigés et les travaux pratiques se font à la maison, hors de
ces 30~heures~: pour les travaux dirigés, la fiche de questions de chaque
chapitre et ses exercices de fin de chapitre (environ dix heures au
total)~; pour les travaux pratiques, les dix laboratoires de NDAP SUKULU,
d'une heure chacun. Les corrigés des exercices sont réunis dans un livret
séparé.

\begin{tableaucours}{cL{6.4cm}cL{3.2cm}}{%
    \tableauHeader{Partie} &
    \tableauHeader{Thème} &
    \tableauHeader{Chap.} &
    \tableauHeader{Niveau}
}
0 & Prologue~: la nuit du 14 mars à Bafoussam & P & introductif \\
I & Fondements~: enquêter sur des traces numériques & 1--2 & fondamental \\
II & Méthode~: acquérir et analyser & 3--7 & fondamental à
intermédiaire \\
III & Cadre normatif et juridique & 8 & fondamental \\
IV & Vulnérabilités~: identifier, coter, corriger & 9--12 & intermédiaire \\
V & Défis avancés & 13--15 & intermédiaire à expert \\
VI & Recherche~: mesurer l'incertitude d'une enquête & 16 & recherche \\
\end{tableaucours}

\section*{Les quinze séances}

\begin{tableaucours}{cL{6.8cm}L{4.2cm}}{%
    \tableauHeader{Séance} &
    \tableauHeader{Contenu} &
    \tableauHeader{Travail à la maison}
}
1 & Prologue et Chap.~1~: l'incident, l'investigation numérique & Fiches P et 1 \\
2 & Chap.~2~: la preuve et la chaîne de possession & Fiche 2, lab 01 \\
3 & Chap.~3~: acquérir et préserver & Fiche 3, lab 02 \\
4 & Chap.~4~: l'analyse du disque & Fiche 4, lab 03 \\
5 & Chap.~5~: l'analyse de la mémoire & Fiche 5, lab 04 \\
6 & Chap.~6~: l'analyse du réseau & Fiche 6, lab 05 \\
7 & Chap.~7~: chronologie et corrélation & Fiche 7, lab 06 \\
8 & Chap.~8~: normes et droit de l'investigation & Fiche 8 \\
9 & Chap.~9~: qu'est-ce qu'une vulnérabilité & Fiche 9, lab 08 \\
10 & Chap.~10~: identifier sans exploiter & Fiche 10, lab 07 \\
11 & Chap.~11~: gérer les vulnérabilités & Fiche 11, lab 09 \\
12 & Chap.~12~: de l'artefact à la vulnérabilité & Fiche 12 \\
13 & Chap.~13~: anti-forensique, nuage et conteneurs & Fiche 13 \\
14 & Chap.~14~: l'enquête intégratrice & Fiche 14, lab 10 \\
15 & Chap.~15 et~16~: le rapport, l'incertitude d'une enquête & Fiches 15
et 16 \\
\end{tableaucours}

Cette répartition est indicative~; l'enseignant l'adapte au calendrier et au
rythme du groupe.

\section*{Le gabarit des chapitres}

Chaque chapitre suit le même plan~: une épigraphe, le niveau et les
prérequis, les objectifs, une situation d'ouverture tirée de l'un des deux
cas, des sections numérotées, ce qu'il faut retenir, des points de
contrôle, trois à quatre exercices gradués, des lectures, et une
transition.

\section*{Les encadrés}

\begin{tableaucours}{L{3.6cm}L{8.6cm}}{%
    \tableauHeader{Encadré} &
    \tableauHeader{Rôle}
}
Objectifs & Ce que l'on doit savoir faire à la fin du chapitre \\
Définition & Les notions, en gras dans le texte et indexées \\
Violation typique & Une erreur fréquente, et la règle qui l'évite \\
Information & Un complément, un renvoi, une nuance \\
Opération (NYANGA, SOLITRANS) & Le cas~: les constatations, puis
l'analyse, avec un niveau de confiance \\
Artefact & Une trace décrite~: source, emplacement, empreinte, ce qu'elle
prouve et ce qu'elle ne prouve pas \\
Chaîne de possession & Ce que le dossier doit contenir pour qu'un autre
puisse refaire l'examen \\
Lecture & Un décompte de cotes \OUI{}~/ \NON{}~/ \NEK{} et sa lecture \\
Métier & Ce que la notion change pour chacun des métiers du cours \\
Texte juridique & Un article cité dans son texte officiel, et sa portée \\
Ce qu'il faut retenir & Les faits et les règles du chapitre \\
Exercice & Trois niveaux de difficulté~: \niveaubase, \niveaumoyen,
\niveauexpert \\
\end{tableaucours}

\section*{Les conventions}

\begin{itemize}
    \item Les cotes~: \OUI{} (établi par une preuve), \NON{} (absent ou
          contredit par une preuve), \NEK{} (connaissance insuffisante pour
          trancher). Une conclusion d'enquête reçoit l'une de ces trois
          cotes.
    \item Constatation et analyse sont toujours séparées~; l'analyse porte
          un niveau de confiance (faible, moyenne, élevée).
    \item Les heures sont données en heure locale du Cameroun (UTC+1), sauf
          mention UTC explicite.
    \item Les textes~: seuls les articles lus dans leur version officielle
          sont cités entre guillemets~; les normes sous licence sont
          paraphrasées~; le degré de vérification figure à
          l'annexe~\ref{ann:textes}.
    \item Les adresses IP des cas appartiennent aux plages réservées à la
          documentation (RFC~5737), et l'identifiant de la faille de la
          passerelle, VULN-2024-4471, est fictif.
    \item Les cas sont fictifs~: noms en majuscules, montants en francs
          CFA, villes réelles.
\end{itemize}
